\documentclass[12pt,a4paper]{article}
\usepackage[utf8]{inputenc}
\usepackage{geometry}
\usepackage{graphicx}
\usepackage{hyperref}
\usepackage{amsmath}
\usepackage{booktabs}
\usepackage{longtable}
\usepackage{enumitem}
\usepackage{parskip}

\geometry{margin=1in}

\title{General Guide to Measurements \\ APMS Aalborg University}
\author{APMS Department}
\date{March 2026}

\begin{document}

\maketitle
\tableofcontents
\newpage

\section*{Author}
I have worked in the field of metrology for more than four years and in measurement for more than ten. I have both received and delivered courses on measurement uncertainty, and have studied several books on the subject, including the recommended book "The Uncertainty in Physical Measurements" by Paolo Fornasini. 

Professionally, I have held roles as a calibration technician and as Technical Manager of an accredited metrology laboratory, conducting both internal and external audits—including a notable ISO/IEC 17025 audit. I have also attended seminars on ISO/IEC 17025 auditor training.
\section{Purpose}
This document provides a comprehensive guide for performing high-quality, traceable, and reproducible measurements at APMS. It is intended to complement the IEEE dataset framework by defining procedures, observation guidelines, and uncertainty evaluation practices.


\vspace{1cm}
\textbf{Rev. 1.0} - Jakob F. Thiesen: Initial release.
\section{Intro}

Metrology — the science of measurements. Most of us spend our entire lives rather oblivious to the fact that a whole branch of science is dedicated to the principle of measurement itself. This, however, is one of the most, if not the most, fundamental branches of science and engineering; without it, neither our theories would be proven nor would our collective and free trade across borders be as smooth as they are today.

Think about this for a moment - have you ever stopped and considered that, over the entire globe, we all agree on the length of 1 cm? This is by no means a coincidence, and furthermore, have you ever thought that 1 cm today is the same as it was 10 years ago? This might sound insane, 1 cm is 1 cm, what is this nutcase on about?

But, alas, this is something that has troubled civilizations for thousands of years. We might make a ruler and declare it the master ruler, make it 1 meter long, and define that this ruler is the definition of 1 m. In fact, this is what we did for many years. Before that, the king's foot was used as a unit of length, as the king was chosen by God, and thus his foot was considered the most accurate — crazy as it might sound. The challenge with a physical ruler, however, is that all materials degrade, and that 1 m might turn out to have changed in 100 years' time.

Armed with this knowledge, we as humans, however, face a problem: what can be trusted? Let me be frank — metrology is, in fact, about questioning everything and quantifying trustworthy numbers that describe how much we do not trust something. This is called \textbf{uncertainty}, and the world is riddled with it. We, as engineers, have a sacred duty to ensure that published data and measurements are of the highest quality, and this falls squarely within the domain of metrology.

Thus, the goal of this document is to introduce engineers to the world of metrology — a field that is both fascinating and, at times, very rigorous. Metrology is the bridge between theory and reality, and an essential tool for every engineer.


\section{Core Concepts in Metrology}

As previously mentioned, metrology is the science of measurement. As will become clear later, a measurement cannot exist without precision, and metrology is therefore closely connected to \textbf{uncertainty}.

To explore the world of metrology, several core concepts must first be defined and described:

\textbf{True Value:}  
This is the value that a physical object would have if it were observed completely correctly, without any measurement error. We may attempt to map or characterize the parameters of an object, but we can never obtain a perfect description. Therefore, the true value is an ideal concept. It exists in reality, since an object has the form and properties that it has, but due to complexity and limitations in our description of the object, we cannot determine the true value exactly.

\textbf{Measurement / Calibration:}  
A procedure and one or more actions that attempt to assess an instrument’s or object’s relationship to one or more true values of a physical quantity, as well as the potential deviation from that true value.

\textbf{Uncertainty:}  
Uncertainty can refer to many things, but as the word itself suggests, in this context it describes the level of confidence associated with a value and its connection to a physical quantity - or rather, the lack of complete certainty we have. In practice, this means a measured value together with a range within which the true value is believed to lie.

\textbf{Standard:}  
A standard (sometimes called a \textit{measurement standard}) is an object used to compare other quantities against. The standard is known with high precision, and therefore its value is trusted to a high degree. It is worth noting that a standard can also be an instrument. In fact, when you take your favorite voltmeter and measure a battery, the voltage measurement you obtain is referenced to the internal voltage reference of the voltmeter.

And thus, from this it should be clear that when we aim to collect data for IEEE, we are in fact probing reality using instruments. There are several best practices to follow to ensure proper measurements are taken. At the core of these are:

\begin{enumerate}
    \item Measurements cannot stand alone; an uncertainty must always accompany the reported value, together with the associated confidence level.
    \item Measurements must be documented and described to such a degree that any other person within your field should be capable of recreating your results, within the bounds of the stated uncertainty - always.
\end{enumerate}

\subsection{Responsibility of the Experimenter}

Up until this point, most students at university have not been given direct responsibility for measurements that may later be used by others. In this context, however, the situation is different.

Metrology relies heavily on a trickle-down effect. Data and measurements are published, reused, and built upon by other researchers and engineers. The measurements you perform today may form the foundation for future research, engineering designs, or automated systems.

In the context of IEEE data collection, the datasets produced may be used to train artificial intelligence systems that can have significant real-world impact. As the person responsible for performing the measurements, you must therefore ensure that the reported measurement values and their associated uncertainties are correct and valid.

If data is published with unrealistically small uncertainties or exaggerated accuracy, systems using that data
- such as AI models - may assign it far more confidence than it deserves. In applications such as autonomous vehicles, drones, or safety-critical systems, this can lead to incorrect decisions with real-world consequences.

Accurate documentation and honest uncertainty estimation are therefore not merely academic requirements; they are part of the professional and ethical responsibility of every engineer performing measurements. Should you later discover that incorrect data has been published, it must be corrected and those using the data must be informed. 

It is important to note that this is not simply a best practice or a suggestion for students - it is an ethical obligation that every engineer carries throughout their professional work.
\section{Scope}

This guide outlines the general principles that should be considered and prepared when experimental data is collected for the purpose of publication. It is intended to be read, understood, and applied by all students and researchers performing experimental measurements in the APMS laboratories, including work related to analog, digital, RF, and AI-based systems.

The main stages covered in this guide are:

\begin{enumerate}
\item \textbf{Preparation - before measurements are performed}
\begin{enumerate}
\item Define the quantity or phenomenon to be measured.
\item Establish the expected results or order of magnitude.
\item Identify the required instruments and equipment.
\item Select and document the measurement method.
\item Define how data will be recorded and documented.
\item Ensure that the required instruments are calibrated or verified against a trusted reference.
\end{enumerate}


\item \textbf{During measurements}
\begin{enumerate}
    \item Record measurement data together with time, operator, instrument identification, and relevant environmental conditions.
    \item Verify instrument settings and monitor measurement stability.
    \item Observe and document unexpected behaviour such as drift, instability, or anomalies.
    \item Repeat measurements as necessary to ensure adequate repeatability.
\end{enumerate}

\item \textbf{After measurements are completed}
\begin{enumerate}
    \item Ensure that all data is stored securely and in an accessible format.
    \item Analyse the data for outliers, anomalies, and unexpected trends.
    \item Perform a general estimation of measurement uncertainty.
    \item Review and update the prepared measurement procedure if errors or improvements are identified.
\end{enumerate}

\end{enumerate}

As mentioned, this document outlines the key areas that must be considered. More detailed, hands-on guidance will be referenced where necessary.

\section{Preparation - Before measurements are performed}
Once a field of study has been defined and experimental data is to be collected, several preparatory aspects should be considered. The experimenter should aim to familiarise themselves with the physical nature of the phenomenon or quantity to be observed and, where possible, develop an intuitive understanding of its expected behaviour. It is acknowledged, however, that complex systems do not always exhibit behaviour that is immediately intuitive.

This section describes the key elements that should be prepared before final measurements are initiated. In practice, this process is often iterative, with preliminary measurements performed to gain familiarity with the instruments, the experimental setup, and the general behaviour of the system under investigation.

\subsection{Define the Quantity}

The first step in any measurement campaign is to clearly define the quantity or phenomenon to be measured or observed. Once this is established, a method for observing or quantifying that quantity must be designed.

In metrology, this is often referred to as an \emph{indirect measurement}, since in many cases we cannot directly observe the quantity of interest. A simple example is measuring electrical current. There are no methods to directly measure current in a straightforward way; instead, we measure a related, observable quantity. Common approaches include:

\begin{itemize}
\item Using a shunt resistor to measure the voltage drop across it.
\item Employing a Hall-effect sensor to measure the magnetic field generated by the current.
\item Measuring the magnetic flux through a coil induced by the current.
\end{itemize}

In each case, we ultimately measure a voltage that is related to the current. This is considered an indirect measurement, where the desired quantity (current) is inferred from a relationship to a directly observable quantity (voltage, magnetic field, or flux).

Although the previous example is simple, the principle of indirect measurement applies broadly. For instance, one might wish to observe the RF spectrum or the reflection coefficients that arise when a hand is moved near an antenna. In this case, the indirect measurement is the reflection coefficient, while the desired phenomenon is the placement or motion of the hand. By measuring the reflection coefficient, we infer the effect of the hand on the antenna, illustrating how indirect measurements allow us to quantify phenomena that cannot be directly observed.

It is important to note that in exploratory work - such as investigating whether the position of a hand can be inferred from measured reflection coefficients - the principle of indirect measurement still applies. In this case, we must not only measure the reflection coefficients, but also know the exact placement of the hand. This allows us to map a relationship between the observable quantity (reflection coefficients) and the underlying phenomenon (hand position), which is essential for drawing meaningful conclusions.

It is also acknowledged that some experimental phenomena are difficult to quantify directly. For example, the observed effect may not be the exact placement of a hand, but rather a hand gesture or motion. In such cases, it may not be meaningful to assign a precise uncertainty to the degree or type of gesture.

The key principle is to remain aware of these limitations and to clearly describe them in the dataset documentation. Even if the phenomenon cannot be rigorously quantified, explicitly stating the context, assumptions, and limitations ensures that subsequent users of the data understand its scope and reliability.

\subsection{Expected Results and Required Instruments}

After the quantity to be measured has been established, the expected range - or at least the order of magnitude - of the quantity should be estimated. This is important for selecting appropriate instruments and settings. In many cases, this process is exploratory, involving small probing measurements to gain insight into the expected behavior. For some measurements, this can be straightforward; for example, measuring a DC voltage with a general-purpose voltmeter is usually simple, as these instruments are well-protected, robust, and reliable.

RF measurements, such as those involving spectrum analyzers or vector network analyzers (VNAs), are considerably more complex. These instruments are expensive and sensitive, and misuse can easily damage them. Understanding the expected magnitude of the signal is therefore critical. Spectrum analyzer measurements, for example, are strongly affected by settings such as resolution bandwidth, attenuator setting, pre-amplifier configuration, and reference level. Incorrect settings can alter the internal measurement chain, causing inaccurate readings, increased intermodulation or distortion, spurious harmonics, or a raised noise floor.

By estimating the expected results in advance, the experimenter ensures that instruments are configured appropriately, measurement quality is optimized, and equipment is protected from potential damage. Admittedly, this can be a tedious task at times. Once the approximate magnitude of the quantity is known, the experimenter should consult the instrument manual to verify proper settings or seek guidance from an experienced technician. However, it is important to do preliminary research before asking for help, in order to use the expert’s time efficiently.

Another important reason to approximate the range of the quantity to be measured is that it allows for a preliminary estimation of measurement uncertainty. This early assessment helps determine whether the chosen method and instruments are suitable for the task at hand, or whether an alternative approach is required. By considering the potential uncertainty upfront, the experimenter can avoid investing time in measurements that are unlikely to yield meaningful or reliable results.

As an extension of this, the expected magnitude and required quality of the measurement should guide instrument selection. For example, highly accurate DC voltage measurements may exceed the capabilities of a standard multimeter. Similarly, for AC voltages, if the frequency is in the hundreds of kilohertz or higher, most general-purpose multimeters will not provide reliable readings. While an oscilloscope is very accurate in the time domain, it is not necessarily precise in amplitude, particularly for small signals.

By considering both the expected magnitude and the required measurement fidelity, the experimenter can choose instruments that are suitable for the task and avoid misinterpretation of results caused by inappropriate equipment. Once again, it is strongly recommended to consult an experienced technician before using any non-general-purpose instruments. These devices are often expensive and sensitive, and improper use can result in damage or inaccurate measurements.
\subsection{Selection and Documentation of Measurement Method}

After the quantity to be measured has been established and the appropriate instruments selected, the experimenter should prepare documentation of the measurement principle. In metrology, this is referred to as a **measurement procedure**, or simply a procedure. This is one of the most important tasks of a measurement engineer: documenting the work in sufficient detail so that it can be reproduced with the same results and the same order-of-magnitude uncertainties.

Preparing a measurement procedure is a rigorous task and requires specifying:

\begin{enumerate}
\item What is being measured.
\item Which instruments are used.
\item How stable the system is (e.g., noise, drift, environmental effects).
\item The estimated measurement uncertainty.
\item How everything is connected (a diagram or picture guide is highly recommended).
\end{enumerate}

This is non-negotiable: if data collection is to be performed, the engineer \textbf{must} present a measurement procedure to both the supervisor and laboratory technicians. This ensures that the correct method is used and prevents damage to instruments or personnel.

It should be noted that in experimental work, the requirements for uncertainty estimation may not be as stringent as in industry. However, the experimenter should always remain aware of potential sources of uncertainty and make rough estimates to ensure that the measured quantity is not overwhelmed or obscured by unknown factors.

The procedure should also describe how data is recorded and documented, whether it is collected manually or through an automated system. If automation is used, the code or scripts responsible for data acquisition should be stored alongside the procedure to ensure reproducibility and traceability.


A separate guide will provide more detailed instructions on designing an appropriate measurement procedure. This current section focuses on the principles and documentation requirements, while the other guide offers practical templates and step-by-step guidance.

It should be noted that the measurement procedure can apply to any system, even something as simple as collecting data using an Arduino and its ADC. In such cases, the procedure should describe how the device is used, how it has been verified or calibrated, and any relevant settings or connections. The purpose of the procedure is to ensure that the measurements can always be reproduced and, if necessary, traced back for verification or analysis.

\subsection{Calibration}

Calibration is a cornerstone of reliable measurements. Instruments used for data collection should always be either **calibrated** or **verified against a trusted, calibrated reference**.

In practice, not all laboratory instruments are externally calibrated, as accredited traceable calibration can be expensive. For example, calibrating a common handheld multimeter can cost more than the instrument itself, while calibration of a vector network analyzer (VNA) can easily range from thousands to tens of thousands of euros, often performed annually.

This underscores the principle that expensive instruments must be handled with care: even a minor accident, such as dropping a VNA calibration kit, can necessitate recalibration, resulting in significant cost.

It is therefore the experimenter’s responsibility to ensure that instruments are either properly calibrated or verified against a known reference. This ensures that the measured data reflect the physical quantity of interest, rather than artifacts arising from inaccurate or malfunctioning instruments.

\section{During Measurements}
After all preparations have been completed, a measurement procedure has been designed, and the required instruments have been selected, the finalized measurement can begin. Ideally, the preparations are so well defined that one simply follows the designed procedure.

However, this is not always the case. This section describes what should be done during measurements and highlights key considerations to be aware of.

\subsection{Records}
Once the measurement begins, the time and place should be recorded, along with the name of the operator, the instruments used, and ideally their serial numbers. Metrology-grade measurements also require recording environmental conditions, such as temperature and humidity. While this may not be strictly necessary for typical experimental work, it is important to be aware of their influence: most laboratory instruments are designed to operate around 23$^\circ$C and 45\% relative humidity.

The environment in which the measurements are performed should always be described. All measurement data should be recorded in accordance with the established measurement procedure.

\subsection{Stability}
During measurements, the experimenter should monitor the system and data for stability, drift, or signs that an instrument may be saturated or operating out of range. If such issues occur, they must be noted, and an evaluation should be made to determine whether the measurements are still acceptable or if the procedure, measurement ranges, or instruments need to be re-assessed. These observations and considerations should be documented and discussed with the supervisor.
\section{After Measurements}
This section describes the considerations that should be made after measurements have been performed. At its core, this involves reflecting on the procedure, what has been measured, and how those measurements compare to what was intended and expected.

It is not uncommon to realize that the measurements may not represent what was originally intended, or that they are only weakly linked to the target or goal. This is also true for exploratory measurements, where reflection may reveal that alternative setups or different quantities would better describe the phenomenon under investigation.


\subsection{Data Storage and Accessibility}
Once all data described in the procedure has been collected, it is imperative that it is stored securely and remains accessible. If cloud services are used, it is good practice to have another person verify access to the data to ensure that it is available and has not been corrupted.

It is important to recognize that high-quality measurement data is highly valuable. For example, a reference laboratory may perform only a limited number of measurements on a 10 V DC standard and issue a signed, accredited certificate at a cost of several thousand euros. 

The collected data should therefore be treated as a valuable asset and stored accordingly, with appropriate backup and access measures in place.

\subsection{Analyze the Data}
The collected data should be carefully analyzed, and the presence of outliers should be evaluated. If outliers are identified, it is essential to determine whether they arise from measurement errors, data entry mistakes, or other identifiable causes.

If an outlier is to be excluded, a clear and well-justified reason must be provided. Ideally, the removal of the outlier should be documented, along with the rationale for its exclusion.

It is important to recognize that the data has been observed under the defined measurement procedure. If the procedure is valid, the origin of an outlier may not be immediately clear-it could be a genuine feature of the phenomenon under investigation, or it could result from misuse of the measurement setup, data entry errors, or other issues.

Data must not be excluded simply because it appears inconsistent or does not fit expectations. Such practice is not scientifically valid, may lead to incorrect conclusions, and is both unethical and in bad faith.

\subsection{General Estimation of Uncertainty}
Depending on the requirements and agreements made with the supervisor, a general 95\% coverage uncertainty (k = 1.96) should be estimated and included in the dataset. It is recognized that, for some scientific work, uncertainty estimation may not always be the primary focus. However, at the core of science lies metrology, and any well-prepared dataset should include some degree of uncertainty estimation.

In many cases, a simplified approach may be sufficient. For example, if measurements are performed using a vector network analyzer under consistent conditions-such as similar attenuator settings, mixer levels, reference levels, and bandwidth-it may be reasonable to assign a uniform uncertainty, e.g., $\pm$ 2 dB, to all measurements.

Such an estimate can be valid, provided it is justified. The basis for the estimation should always be clearly described, even if a simplified method is used. I would like to quote the international recognized bible (BIPM GUM) within uncertainties here:

\begin{quote}
\centering
\textbf{In general, the result of a measurement is only an approximation or estimate of the value of the measurand and is complete only when accompanied by a statement of the uncertainty of that estimate.}
\end{quote}

Please be aware that the evaluation of uncertainty is not something that should be delegated to technicians. While they may assist with instrument setup and configuration, the evaluation and estimation of uncertainty is the responsibility of the experimenter and the designer of the measurement procedure.

There are numerous documents and general guides available to support this process, several of which are referenced throughout this document.


\subsection{Review and Update Measurement Procedure}
The final step is to review the measurement procedure and evaluate whether any elements should be removed, added, or modified. The procedure must always describe the exact method by which the data was collected.

If the procedure is changed and the collected data no longer aligns with the documented method, the data is no longer valid under that procedure. In such cases, the data must either be collected again or documented under a separate procedure specific to that dataset.

There is no wiggle room in this principle: data must always follow the documented method. This is the cornerstone of reliable and traceable measurements, ensuring that every dataset can be understood, reproduced, and validated.

It must be made absolutely clear: under no circumstances should data or procedures be altered to make results appear more favorable. Any attempt to manipulate measurements undermines the scientific method and the reliability of the dataset. The integrity of the measurement process is paramount; it is not subject to personal preference, intuition, or aesthetic judgment. 

\subsection{Publication}
Once all the preceding steps have been completed and both the experimenter and supervisor are satisfied, the data is ready for publication. If this document has been followed thoroughly, the resulting dataset should be of such quality that it rivals that of the finest laboratories in industry. 

High-quality measurements are a source of pride, and producing them requires effort, diligence, and attention to detail. Following rigorous procedures ensures that the data is trustworthy, reproducible, and valuable to others in the scientific and engineering community.


\section{References}
\begin{itemize}[noitemsep]
    \item \textbf{ISO/IEC 17025:} \textit{General requirements for the competence of testing and calibration laboratories}
    \item \textbf{BIPM GUM:} \textit{Guide to the Expression of Uncertainty in Measurement}
    \item \textbf{DANAK AB11:} Measurement uncertainty in calibration
    \item \textbf{DANAK AB12:} Evaluation and reporting of measurement uncertainty
\end{itemize}



\section{Roles and Responsibilities}

\subsection{Student}
The student is responsible for performing the measurements, developing the measurement procedure, logging all data, and observing any anomalies. In general, the student ensures that the measurements are conducted correctly, the data is accurately collected, and all records are properly stored.

\subsection{Supervisor}
The supervisor provides guidance to the student as needed. This may include assistance with measurement setup, understanding the underlying principles, determining what to measure, and how to describe or document the measurements. The supervisor also reviews the procedures and ensures that the resulting data is reliable and reproducible.

\section{Miscellaneous}
As a final note, while all of this may initially seem “a bit much,” once you become familiar with the world of Metrology, it is both exciting and fascinating. The field is filled with impressive discoveries and achievements. One example that particularly fascinates me is the Kibble balance, which links the kilogram to Planck's constant through the electromagnetic force -a truly remarkable accomplishment.

The following section contains references to guides on uncertainty calculations, measurement procedures, and other resources for further study.

\end{document}